% =====================================================================
%  code.tex — Coloration syntaxique (listings), pensée pour le noir et blanc
%
%  Styles disponibles : python, javascript, nodejs, html, css, java,
%                       dockerfile, yaml
%
%  Utilisation dans une question :
%
%    \begin{code}{python}              % bloc de code
%    print("Bonjour")
%    \end{code}
%
%    \begin{code}[numbers=none]{css}   % options listings entre crochets
%    ...
%    \end{code}
%
%    \codefile{java}{questions/code/Compte.java}      % depuis un fichier
%    \codefile[firstline=3, lastline=10]{java}{...}   % extrait de lignes
%
%    \codeinline{x = 3}               % code dans le texte
%
%  Règles d'impression : mots-clés en gras, commentaires en italique gris,
%  chaînes en gris foncé.
% =====================================================================

% --- Teintes de gris (noms requis : « ! » pose problème dans listings) -
\colorlet{codecomment}{black!60}
\colorlet{codestring}{black!75}
\colorlet{codenumber}{black!55}
\colorlet{codeframe}{black!45}
\colorlet{codearrow}{black!50}
\colorlet{codeattribute}{black!70}

% --- Style commun -----------------------------------------------------
\lstdefinestyle{base}{
  basicstyle        = \ttfamily\small,
  keywordstyle      = \bfseries,
  commentstyle      = \itshape\color{codecomment},
  stringstyle       = \color{codestring},
  numbers           = left,
  numberstyle       = \scriptsize\color{codenumber},
  numbersep         = 8pt,
  xleftmargin       = 2.2em,
  framexleftmargin  = 0.3em,
  frame             = single,
  framerule         = 0.4pt,
  rulecolor         = \color{codeframe},
  columns           = fullflexible,
  keepspaces        = true,
  showstringspaces  = false,
  tabsize           = 4,
  upquote           = true,
  extendedchars     = true,
  breaklines        = true,
  postbreak         = \mbox{\textcolor{codearrow}{$\hookrightarrow$}\space},
  aboveskip         = 0.6em,
  belowskip         = 0.6em,
}

% --- Définitions de langages absents de listings ---------------------
\lstdefinelanguage{JavaScript}{
  morekeywords = {async, await, break, case, catch, class, const, continue,
    default, delete, do, else, export, extends, finally, for, from,
    function, if, import, in, instanceof, let, new, of, return, static,
    super, switch, this, throw, try, typeof, var, void, while, yield,
    true, false, null, undefined},
  sensitive    = true,
  morecomment  = [l]{//},
  morecomment  = [s]{/*}{*/},
  morestring   = [b]',
  morestring   = [b]",
  morestring   = [b]`,
}

\lstdefinelanguage{CSS}{
  alsoletter   = {-@!},
  morekeywords = {color, background, background-color, background-image,
    border, border-radius, border-color, border-width, border-style,
    margin, margin-top, margin-right, margin-bottom, margin-left,
    padding, padding-top, padding-right, padding-bottom, padding-left,
    width, height, max-width, min-width, max-height, min-height,
    display, position, top, right, bottom, left, z-index, float, clear,
    font, font-family, font-size, font-weight, font-style, line-height,
    text-align, text-decoration, text-transform, letter-spacing,
    flex, flex-direction, flex-wrap, justify-content, align-items, gap,
    grid, grid-template-columns, grid-template-rows, grid-area,
    opacity, overflow, visibility, cursor, transition, transform,
    animation, box-shadow, box-sizing, content, list-style},
  morekeywords = [2]{@media, @import, @keyframes, @font-face, !important},
  sensitive    = false,
  morecomment  = [s]{/*}{*/},
  morestring   = [b]",
  morestring   = [b]',
}

\lstdefinelanguage{HTML5}[]{HTML}{
  morekeywords = {header, nav, main, section, article, aside, footer,
    figure, figcaption, video, audio, source, canvas, template,
    button, label, select, option, textarea, output, progress,
    charset, lang, placeholder, required, defer, async, type,
    for, data, alt, src, href, rel, id, class, name, value},
  sensitive    = false,
}

% Instructions en majuscules (convention) : « from » dans un RUN reste normal
\lstdefinelanguage{Dockerfile}{
  morekeywords = {FROM, AS, RUN, CMD, ENTRYPOINT, COPY, ADD, WORKDIR, ENV,
    ARG, EXPOSE, VOLUME, USER, LABEL, HEALTHCHECK, SHELL, STOPSIGNAL,
    ONBUILD},
  sensitive    = true,
  morecomment  = [l]{\#},
  morestring   = [b]",
  morestring   = [b]',
}

% Clés en gras : tout est gras jusqu'au premier « : » de la ligne,
% puis le reste de la ligne (la valeur) revient au style normal
\lstdefinelanguage{YAML}{
  morekeywords = {true, false, null, yes, no, on, off},
  sensitive    = false,
  morecomment  = [l]{\#},
  morestring   = [b]",
  morestring   = [b]',
  moredelim    = **[l][\mdseries]{:},
}

% --- Styles par langage ----------------------------------------------
\lstdefinestyle{python}{
  style        = base,
  language     = Python,
  morekeywords = {True, False, None, async, await, nonlocal, with, as,
                  assert, yield, lambda},
}

\lstdefinestyle{javascript}{
  style    = base,
  language = JavaScript,
}

% Node.js = JavaScript + quelques identifiants propres à Node
\lstdefinestyle{nodejs}{
  style        = javascript,
  morekeywords = {require, module, exports, process},
}

\lstdefinestyle{html}{
  style            = base,
  language         = HTML5,
  tagstyle         = \bfseries,
  keywordstyle     = \color{codeattribute},
  markfirstintag   = true,
  usekeywordsintag = true,
}

\lstdefinestyle{css}{
  style         = base,
  language      = CSS,
  keywordstyle  = \bfseries,
  keywordstyle  = [2]\bfseries\itshape,
}

\lstdefinestyle{java}{
  style        = base,
  language     = Java,
  morekeywords = {var, record, sealed, permits, yield},
}

\lstdefinestyle{dockerfile}{
  style    = base,
  language = Dockerfile,
}

\lstdefinestyle{yaml}{
  style        = base,
  language     = YAML,
  basicstyle   = \ttfamily\small\bfseries,
  keywordstyle = \mdseries\itshape,
  commentstyle = \mdseries\itshape\color{codecomment},
}

\lstset{style=base}

% --- Commandes pour les questions ------------------------------------
% Bloc de code :  \begin{code}[options]{langage} ... \end{code}
\lstnewenvironment{code}[2][]{\lstset{style=#2, #1}}{}

% Code depuis un fichier : \codefile[options]{langage}{chemin}
\NewDocumentCommand{\codefile}{O{} m m}{%
  \lstinputlisting[style=#2, #1]{#3}%
}

% Code dans le texte : \codeinline{...} ou \codeinline|...|
% (ne pas l'utiliser dans l'argument d'une autre commande)
\NewDocumentCommand{\codeinline}{v}{\texttt{#1}}
